\chapter{相对运动}\label{ch:g10:relative-motion}

一名乘客手拿咖啡，以悠闲的步行速度走向高速列车车头。
站在铁轨旁的观察者却测得，此人的速度超过 \qty{300}{km/h}。
两种测量结果都正确。
两者都不是乘客的“真实”速度，因为根本不存在这样的速度：运动并非物体单独具有的性质，
而是物体\emph{与}描述它时所选取的观察视角共同决定的。
本章将精确刻画这一视角 ---
并教你如何选取一个好的视角。

\section{运动是相对的}

\begin{definition}[参考系]\label{def:g10:relative-motion:frame}
\emph{参考系}\index{参考系}是一个刚体 --- 例如地面、列车车厢或整个地球 ---
相对于它记录位置，并配有一只用来记录时刻的时钟。
若物体在某参考系中的位置不随时间改变，
则称该物体在此参考系中\emph{静止}\index{静止}；
否则称其在此参考系中运动。
\end{definition}

\begin{definition}[轨迹]\label{def:g10:relative-motion:trajectory}
物体在给定参考系中的\emph{轨迹}\index{轨迹}，
是它在该参考系中依次经过的各个位置所连成的曲线：
可以是直线、圆周，也可以是一段弧线……
只有先指明参考系，轨迹才有意义。
\end{definition}

\begin{example}[列车上的乘客]\label{ex:g10:relative-motion:walker}
列车相对地面以 \qty{300}{km/h} 的速度行驶，
同时乘客相对列车以 \qty{4}{km/h} 的速度向车头走去。
\begin{itemize}
  \item 在\emph{列车}参考系中：乘客沿过道以 \qty{4}{km/h} 的速率前进；
        座位静止；窗外的树木向后飞驰，
        速率为 \qty{300}{km/h}。
  \item 在\emph{地面}参考系中：座位以 \qty{300}{km/h} 运动；
        乘客的速度为 \qty{304}{km/h}，树木静止。
\end{itemize}
同一场景，两个参考系，两种截然不同的描述 ---
而在车厢内进行的任何实验，
都无法断定其中哪一种才是“真实的”。
\end{example}

\begin{proposition}[运动的相对性]\label{prop:g10:relative-motion:relativity}
物体的静止或运动状态、轨迹及速率，
都取决于描述它们时所选取的参考系。
\end{proposition}

\begin{proof}
一个例子即可说明每一项结论。
由\cref{ex:g10:relative-motion:walker}可知，坐着的乘客在列车参考系中静止，
而在地面参考系中以 \qty{300}{km/h} 的速度运动（静止与速率都取决于参考系）。
至于轨迹，可观察自行车车轮上的气门嘴：骑车人看到它绕轮轴画出一个圆周，
而站在路边的观察者看到它描出一连串拱形曲线，
这种曲线称为\emph{摆线}\index{摆线}，
见下图。
\end{proof}

\begin{omfigure}
\begin{tikzpicture}[scale=0.85, font=\small]
  \draw[gray!60] (-1.6,0) -- (1.6,0);
  \draw[gray, thin] (0,1) circle (1);
  \fill (0,1) circle (1.5pt);
  \draw[gray] (0,1) -- (0.71,1.71);
  \draw[omDef, thick, -Stealth] (1,1) arc (0:330:1);
  \fill[omThm] (0.71,1.71) circle (2pt) node[above right] {气门嘴};
  \node[omDef] at (0,-0.5) {自行车车架：圆周};
\end{tikzpicture}
\qquad
\begin{tikzpicture}[scale=0.42, font=\small]
  \draw[gray!60] (-0.6,0) -- (13.2,0);
  \draw[omDef, thick] plot[domain=0:12.566, samples=100, smooth]
    ({\x - sin(\x r)}, {1 - cos(\x r)});
  \draw[gray, thin] (9.42,1) circle (1);
  \fill (9.42,1) circle (1.5pt);
  \draw[gray] (9.42,1) -- (9.42,2);
  \fill[omThm] (9.42,2) circle (3pt);
  \node[omDef] at (6.3,-1.0) {路面参考系：摆线};
\end{tikzpicture}

{\small 同一气门嘴，两种参考系：对骑车人是圆周，
对路边观察者是摆线（一连串拱形）。}
\end{omfigure}

\begin{remark}
注意两幅图中\emph{不变}的是什么：
气门嘴、车轮与物理规律都未改变，变化的只是描述。
问“气门嘴是否真的在做圆周运动”，
就如同问山是否“真的”在左边：这取决于你站在哪里。
\end{remark}

\section{平均速率}

\begin{definition}[平均速率]\label{def:g10:relative-motion:speed}
物体在给定参考系中的\emph{平均速率}\index{速率!平均}，
是它在该参考系中经过的路程 $d$ 与
所用时间 $t$ 之比：
\[
v = \frac{d}{t} .
\]
若 $d$ 以米为单位、$t$ 以秒为单位，则 $v$ 以米每秒
（\unit{m/s}）为单位；日常生活中更常用千米每小时（\unit{km/h}）。
\end{definition}

\begin{method}[米每秒与千米每小时之间的换算]\label{met:g10:relative-motion:convert}
一千米为 \qty{1000}{m}，一小时为 \qty{3600}{s}，因此
$\qty{1}{m/s} = \qty{3.6}{km/h}$：
\begin{itemize}
  \item 由 \unit{m/s} 换算为 \unit{km/h}：\emph{乘以} $3.6$；
  \item 由 \unit{km/h} 换算为 \unit{m/s}：\emph{除以} $3.6$。
\end{itemize}
验算方法：\unit{km/h} 的数值总是两者中
\emph{较大}的那一个。
\end{method}

\begin{example}[几个值得了解的速率]\label{ex:g10:relative-motion:convert}
一次轻快的步行：$\qty{1.5}{m/s} \times 3.6 = \qty{5.4}{km/h}$。一辆城市汽车
以 \qty{54}{km/h} 行驶：$54 / 3.6 = \qty{15}{m/s}$ --- 相当于
每秒约行驶 \qty{15}{m}，这一数值在走下路缘之前值得记住。
空气中的声音：$\qty{340}{m/s} \times 3.6 = \qty{1224}{km/h}$。
\end{example}

\begin{example}[中途休息的徒步]\label{ex:g10:relative-motion:hike}
一名徒步者步行 \qty{2}{h} 行进 \qty{12}{km}，然后休息
\qty{1}{h}，随后在最后 \qty{1.5}{h} 内行进 \qty{6}{km}。全程
的平均速率：
\[
v = \frac{d}{t} = \frac{12 + 6}{2 + 1 + 1.5}
= \frac{\qty{18}{km}}{\qty{4.5}{h}} = \qty{4.0}{km/h}.
\]
平均速率计入\emph{全部}经过的时间，包括休息时间 --- 它一般
不等于各分段速率的平均值。
\end{example}

\section{选择参考系：地面、地心、日心}

由于每个参考系都合理，应选择使所关心的运动
看起来最简单的参考系。三种
常用选择即可涵盖本书的大部分物理内容。

\begin{definition}[三种标准参考系]\label{def:g10:relative-motion:frames}
\begin{itemize}
  \item \emph{地面参考系}\index{地面参考系}（或
        \emph{实验室参考系}\index{实验室参考系}）固连于
        地球表面进行实验之处。它是
        日常运动的适用参考系：自由下落的物体、车辆、
        体育运动和实验台上的物体。
  \item \emph{地心参考系}\index{地心参考系}以地球为中心，
        但\emph{不}随地球一起自转（其方向
        指向遥远的恒星）。它是月球和
        人造卫星的适用参考系，二者整体绕地球运行并
        忽略地球的每日自转。
  \item \emph{日心参考系}\index{日心参考系}以太阳为中心，
        其方向同样固定于遥远的
        恒星。它是行星的适用参考系：每颗行星随后沿
        近似圆形的简单轨道运动。
\end{itemize}
\end{definition}

\begin{omfigure}
\begin{tikzpicture}[scale=1.0, font=\small]
  \fill[omProp] (0,0) circle (0.26);
  \node[below=3pt] at (0,-0.26) {太阳};
  \draw[dashed, gray] (0,0) circle (2.5);
  \draw[-Stealth, gray] (2.5,0) arc (0:22:2.5);
  \fill[omDef] (2.5,0) circle (0.12);
  \node[below right] at (2.55,-0.10) {地球};
  \draw[dashed, gray] (2.5,0) circle (0.75);
  \draw[-Stealth, gray] (3.25,0) arc (0:55:0.75);
  \fill[gray] (2.5,0.75) circle (0.06);
  \node[above] at (2.5,0.80) {月球};
  \draw[omThm] (2.38,0.12) -- (1.2,1.35);
  \node[omThm, left, align=right] at (1.25,1.45)
    {地面参考系：\\随地表一起转动的点};
  \node[omDef, align=left, anchor=west] at (3.6,0.1)
    {地心参考系：\\以地球为中心\\(月球、卫星)};
  \node[omProp, align=left, anchor=west] at (-1.4,-1.6)
    {日心参考系：以太阳为中心（行星）};
\end{tikzpicture}

{\small 三个嵌套的参考系：地面随地球一起转动，地心
参考系随地球一起运动而不自转，日心
参考系从太阳处俯瞰整个运动过程。}
\end{omfigure}

\begin{method}[选择参考系]\label{met:g10:relative-motion:choose}
先指明需要描述运动的物体，再选取
使该运动最简单的参考系：
\begin{enumerate}
  \item 地表附近、持续数分钟的运动 --- 地面
        参考系；
  \item \emph{绕}地球运行的天体（月球、卫星）--- 地心
        参考系；
  \item 绕太阳运行的天体（行星、彗星）--- 日心参考系。
\end{enumerate}
描述运动之前，务必先说明所选参考系：未注明参考系的轨迹或
速率毫无意义
（见\cref{prop:g10:relative-motion:relativity}）。
\end{method}

\begin{remark}[谁是对的：托勒密还是哥白尼？]\label{rem:g10:relative-motion:history}
“太阳从东方升起”与“地球向东自转”
是同一观测结果在两个参考系中的表述：前者在地面参考系中
描述太阳划过天空，后者在地心参考系中
描述观测者随地球绕地轴转动。几个
世纪以来，参考系本身就是争论的战场：托勒密的地心
体系描述精确地追踪了天空，但为追踪行星需要
层层叠叠的环路才能跟随行星，而
哥白尼于 1543 年复兴的日心描述使每颗行星都成为简单的近圆轨道，
并把火星的奇异逆行回路解释为仅仅是参考系效应 ---
地球超越了较慢的行星。现代的结论并非一个参考系为真、另一个为假，
而是两者并不同样
\emph{方便}；选择某些参考系（即惯性参考系）的
更深层准则，将随惯性原理一并
在\cref{ch:g10:inertia}中给出。
\end{remark}

\section{沿直线的运动合成}

步行者在火车参考系中的 \qty{4}{km/h}，与
地面参考系中的 \qty{304}{km/h}，二者有何联系？沿同一直线，
答案很简单：直接相加。

\begin{proposition}[一维速率的合成]\label{prop:g10:relative-motion:compose}
设一物体在沿直线运动的支撑物（如步道、
火车、流动的水）上运动，支撑物本身沿同一直线
相对地面滑动。取直线上的一个正方向，
并用正负号计入速率，
\[
v_{\text{物体/地面}} \;=\; v_{\text{物体/支撑物}} \;+\;
v_{\text{支撑物/地面}} .
\]
用文字表述：同向的速率相加；反向的
速率相减。
\end{proposition}

\begin{proof}
在一秒内，支持物带着物体沿直线移动 $v_{\text{支持物/地面}}$
米，物体又沿支持物前进 $v_{\text{物体/支持物}}$
米。两段位移都在同一直线上，
因此物体在这一秒内相对地面的总位移
为二者的代数和 --- 即其相对地面的速率。
\end{proof}

\begin{example}[移动人行道]\label{ex:g10:relative-motion:walkway}
机场人行道以 \qty{1.0}{m/s} 的速率运行。
一名旅客顺着运行方向，相对传送带以 \qty{1.5}{m/s} 的速率行走：
相对地面的速率为 $1.5 + 1.0 = \qty{2.5}{m/s}$。若逆向以同样速率行走，
则为 $1.5 - 1.0 = \qty{0.5}{m/s}$，仍（缓慢地）
逆着传送带前进。若站在传送带上不动：
相对地面的速率为 \qty{1.0}{m/s}，相对
传送带的速率为 \qty{0}{m/s} --- 同一时刻，人在两个不同的
参考系中既静止又运动。
\end{example}

\begin{omfigure}
\begin{tikzpicture}[font=\small]
  \draw[fill=gray!12] (0,0) rectangle (7,0.3);
  \node[gray] at (6.3,0.15) {传送带};
  \draw (2.5,0.3) -- (2.5,1.05);
  \draw (2.5,0.75) -- (2.2,0.45) (2.5,0.75) -- (2.8,0.45);
  \fill (2.5,1.2) circle (2.5pt);
  \draw[-Stealth, omDef, thick] (0.4,-0.45) -- (1.9,-0.45)
    node[midway, below] {传送带/地面：\qty{1.0}{m/s}};
  \draw[-Stealth, omThm, thick] (2.5,1.6) -- (4.75,1.6)
    node[midway, above] {行人/传送带：\qty{1.5}{m/s}};
  \draw[-Stealth, omMeth, very thick] (2.5,2.3) -- (6.25,2.3)
    node[midway, above] {行人/地面：\qty{2.5}{m/s}};
\end{tikzpicture}

{\small 一维合成：箭头按比例绘制，
地面参考系中的箭头为另外两个箭头之和。}
\end{omfigure}

\begin{remark}[两个方向同时存在]\label{rem:g10:relative-motion:twodim}
当两个运动\emph{不}在同一直线上时——例如船头
正对河岸航行，而水流把船冲向下游——
两个运动各自独立进行，速度按
首尾相接的箭头合成，与 One Math Book 中的矢量相同。
船仍只在由其横渡速率决定的时间内到达对岸，
但会在瞄准点的下游靠岸，并且
其在地面参考系中的轨迹为一条倾斜的直线。完整的矢量
处理方法将在后面的章节中介绍，这里只是读图。
\end{remark}

\begin{omfigure}
\begin{tikzpicture}[scale=0.95, font=\small]
  \fill[omDef!10] (0,0) rectangle (6.4,2.4);
  \draw (0,0) -- (6.4,0) (0,2.4) -- (6.4,2.4);
  \draw[-Stealth, gray] (4.8,0.35) -- (5.6,0.35);
  \node[gray] at (5.2,0.62) {水流};
  \draw[dashed, omMeth] (1.5,0) -- (3.3,2.4);
  \draw[-Stealth, omThm, thick] (1.5,0.6) -- (1.5,2.0)
    node[left, pos=0.7] {船/水};
  \draw[-Stealth, omDef, thick] (1.5,0.6) -- (2.55,0.6)
    node[below, pos=0.8] {水/地面};
  \draw[-Stealth, omMeth, very thick] (1.5,0.6) -- (2.55,2.0)
    node[right, pos=0.75] {船/地面};
  \fill (1.5,0.6) circle (1.8pt);
\end{tikzpicture}

{\small 船头正对河岸航行：在水参考系中，船
沿直线横渡；在地面参考系中，船沿倾斜的虚线
行进，横渡过程中向下游漂移。}
\end{omfigure}

\section{练习}

\begin{exercise}[$\star$]\label{exo:g10:relative-motion:1}
一名乘客在以 \qty{100}{km/h} 行驶的火车上睡着了。请判断
该乘客相对于下列对象是静止还是运动，若运动，速率为多少：
车厢；地面；铁轨旁的一棵树；
迎面驶来、速率同为 \qty{100}{km/h} 的另一列火车。然后
描述在乘客所乘火车的参考系中，这棵树的运动。
\end{exercise}

\begin{exercise}[$\star$]\label{exo:g10:relative-motion:2}
换算：将 \qty{72}{km/h} 和 \qty{108}{km/h} 化为 \unit{m/s}；
将 \qty{25}{m/s} 和 \qty{340}{m/s} 化为 \unit{km/h}。检验每个答案时，
请运用\cref{met:g10:relative-motion:convert}中的合理性检验法则。
\end{exercise}

\begin{exercise}[$\star$]\label{exo:g10:relative-motion:3}
一名骑车人用时 \qty{1}{h}~\qty{30}{min} 行驶了 \qty{36}{km}；一名短跑运动员
用时 \qty{10.0}{s} 跑完了 \qty{100}{m}。计算两人的平均速率，分别以
\unit{m/s} 和 \unit{km/h} 为单位。谁更快？这个结果是否
出乎意料？
\end{exercise}

\begin{exercise}[$\star$]\label{exo:g10:relative-motion:4}
一条移动人行道以 \qty{0.9}{m/s} 的速率运行。一名旅客以
相对传送带 \qty{1.4}{m/s} 的速率行走。求该旅客相对地面的速率，
分别在顺着传送带运行方向行走、逆着运行方向行走，以及
站在传送带上不动时的情况。
\end{exercise}

\begin{exercise}[$\star$]\label{exo:g10:relative-motion:5}
对每种运动，从地面、地心和日心参考系中选出最方便的一个，
并用一句话说明理由：一支笔在行驶的火车内
掉落（附加题：还有一个更好的参考系）；国际空间站的
轨道；一个火星年；一次
\qty{100}{m} 短跑。
\end{exercise}

\begin{exercise}[$\star\star$]\label{exo:g10:relative-motion:6}
一列火车以 \qty{108}{km/h} 的速率行驶。一名乘客相对火车，以
\qty{1.5}{m/s} 的速率向车头方向行走。
\begin{enumerate}
  \item 将火车速率换算为 \unit{m/s}，然后求
        乘客相对地面的速率，分别为向车头和
        向车尾行走时的速率。
  \item 车厢长 \qty{60}{m}。乘客走到车头需要多长时间，
        以及在此期间乘客相对
        地面行进了多远？
\end{enumerate}
\end{exercise}

\begin{exercise}[$\star\star$]\label{exo:g10:relative-motion:7}
一艘船在静水中的速率为 \qty{5.0}{m/s}；河水以
\qty{2.0}{m/s} 的速率流动。船顺流行驶 \qty{420}{m}，掉头后
返回。
\begin{enumerate}
  \item 求船在每段行程中相对河岸的速率，以及
        每段行程的时长。
  \item 求往返全程的平均速率。为什么它小于
        \qty{5.0}{m/s}，尽管水流在一段行程中助力，
        在另一段行程中阻碍的程度恰好相同？
\end{enumerate}
\end{exercise}

\begin{exercise}[$\star\star$]\label{exo:g10:relative-motion:8}
两列火车在平行轨道上错车而过：列车 A
（长 \qty{180}{m}）以 \qty{25}{m/s} 的速率行驶，列车 B（长 \qty{90}{m}）以 \qty{20}{m/s} 的速率行驶。
\begin{enumerate}
  \item 两车反向行驶。在 A 的参考系中，B 的速率为多少，
        以及完整的错车过程持续多久
        （从车头相遇到车尾分开）？
  \item 若两车同向行驶，上述问题的答案如何？
\end{enumerate}
\end{exercise}

\begin{exercise}[$\star\star$]\label{exo:g10:relative-motion:9}
一辆自行车沿直线道路行驶。描述车胎气门嘴在下列参考系中的轨迹：
骑车人的参考系、地面参考系，以及气门嘴自身的参考系。
另外，描述车轮轮毂在地面参考系中的轨迹。
无需计算 --- 说出每条曲线的名称。
\end{exercise}

\begin{exercise}[$\star\star$]\label{exo:g10:relative-motion:10}
地球赤道长约 \qty{40000}{km}，地球自转
一周需要 \qty{24}{h}。
\begin{enumerate}
  \item 求赤道上一点在地心参考系中的速率，
        分别以 \unit{km/h} 和 \unit{m/s} 表示。
  \item 各用一句话解释：为什么我们感觉不到这一速率，以及
        为什么“太阳升起”和“地球自转”描述的是同一个
        观察结果。
\end{enumerate}
\end{exercise}

\begin{exercise}[$\star\star$]\label{exo:g10:relative-motion:11}
站在运行的自动扶梯上静止不动，上行需时 \qty{60}{s}。
若扶梯停运，沿同一扶梯步行上行需时 \qty{90}{s}。
若在运行的扶梯上步行上行，需时多久？（请用速率而不是时间
进行计算。）
\end{exercise}

\begin{exercise}[$\star\star\star$]\label{exo:g10:relative-motion:12}
一名司机两次行驶相同的距离：去程速率为 \qty{60}{km/h}，
回程速率为 \qty{90}{km/h}。
\begin{enumerate}
  \item 证明往返全程的平均速率为
        \qty{72}{km/h} --- 而非 \qty{75}{km/h}。
  \item 解释直觉上的答案错在哪里，并指出哪一段行程
        对平均速率起主导作用。
\end{enumerate}
\end{exercise}

\begin{exercise}[$\star\star\star$]\label{exo:g10:relative-motion:13}
两名骑车人相距 \qty{35}{km}，相向而行，一人速率为
\qty{15}{km/h}，另一人为 \qty{20}{km/h}。
\begin{enumerate}
  \item 在第一名骑车人的参考系中，描述第二名骑车人的运动，
        并推断两人何时相遇。
  \item 两人相遇时，各自骑行了多远？
  \item 若两人改为同向骑行（速率较快者在后），
        则追及过程持续多久？
\end{enumerate}
\end{exercise}

\begin{exercise}[$\star\star\star$]\label{exo:g10:relative-motion:14}
一名游泳者横渡宽 \qty{60}{m} 的河流，相对于水以
\qty{1.2}{m/s} 的速率垂直对岸游去；水流速率为
\qty{0.9}{m/s}。
\begin{enumerate}
  \item 解释为什么渡河时间仅由游泳者的横渡速率决定，
        并求出渡河时间。
  \item 游泳者在下游多远处靠岸？
  \item 利用直角三角形，求游泳者在地面参考系中的速率，
        以及其实际在地面上游过的路径长度。
\end{enumerate}
\end{exercise}

\begin{exercise}[$\star\star\star$]\label{exo:g10:relative-motion:15}
地球绕太阳运行，轨道近似为一个圆，半径
为 \qty{1.50e11}{m}，绕行一周需一年（$\num{3.16e7}$ 秒）。
\begin{enumerate}
  \item 求地球在日心参考系中的速率，以
        \unit{km/s} 表示。
  \item 将其与在
        \cref{exo:g10:relative-motion:10}中求得的赤道自转速率进行比较。
  \item 从地球上看，火星有时似乎会停下来，
        在恒星之间向后漂移数周。请用两句话，借助本章词汇，
        说明为何日心参考系能够化解这一
        谜团。
\end{enumerate}
\end{exercise}

\section{问题：移动人行道、河流与卫星}

\begin{problem}[{周末问题 --- 移动人行道、河流与卫星：谁在运动，
相对于什么运动，以及“安乐椅旅行者”的
最终速率报告}]
\label{pb:g10:relative-motion:1}
机场、渡船过河以及黄昏天空中缓缓滑过的一个亮点：
三个日常场景，同一个问题问了三遍 --- 谁在运动，相对于什么运动？
每一部分都选定一个参考系、进行计算，并观察答案如何随观察点的不同而改变；
结尾一问汇总的是此刻 \emph{你} 的运动速率——你端坐不动，
却究竟在以多大的速率运动？

\medskip
=\noindent\textbf{第一部分 --- 移动人行道赛跑。}
机场人行道长 \qty{90}{m}，以 \qty{1.0}{m/s} 的速率运行。一名
旅客在其上行走，速率为 \qty{1.5}{m/s}（相对于脚下的人行道）。
=\begin{enumerate}
  \item 人行道把静止站立的旅客
        从一端运到另一端需要多长时间？
  \item 旅客在运动的传送带上、与其运动方向相同地行走时，
        需要多长时间？
  \item 旅客在静止的地面上、紧贴人行道行走时，需要多长时间？
  \item 一名恶作剧者在传送带上逆向行走，速率为 \qty{1.5}{m/s}
        （相对于传送带）。求其相对于地面的速率，以及他从远端
        回到入口所需的时间？
  \item 一个孩子以恰好 \qty{1.0}{m/s} 的速率逆向奔跑
        （相对于传送带）。在地面参考系中描述该孩子的运动，
        并指出这相当于重新发明了哪种健身器械。
  \item 用两列表格汇总第 1--5 题的结果：每个运动者
        列出其在传送带参考系和地面参考系中的速率。旅客的腿部
        感受到的是哪一列的速率？
=\end{enumerate}

=\medskip
=\noindent\textbf{第二部分 --- 河流。}
一艘渡船横渡一条宽 \qty{120}{m} 的河流。其相对于
水的速率为 \qty{2.0}{m/s}，水流速率为 \qty{1.5}{m/s}。驾驶员
首先将船头直指对岸。
=\begin{enumerate}[resume]
  \item 渡河需要多长时间？
  \item 渡船在下游多远处靠岸？
  \item 利用直角三角形，求渡船在地面参考系中的速率，
        以及其在地面参考系中轨迹的长度。
        （会出现一组熟悉的勾股数。）
  \item 分别在水参考系和地面参考系中描述渡船的轨迹，
        二者有何相同，又有何不同？
  \item 一名游泳者相对于水的速率为 \qty{1.0}{m/s}，
        从同一岸出发，像渡船那样径直朝对岸
        游去。他渡河需要多长时间？在下游多远处
        靠岸？求其在地面参考系中的速率
        （仍用直角三角形）。为什么速度较慢的游泳者
        比渡船漂移得更远？
  \item 一根木头从渡船旁漂过。给出其在地面参考系中的速率
        和在水参考系中的速率。木头在哪一参考系中“在运动”？
        请加以评述。
=\end{enumerate}

=\medskip
=\noindent\textbf{第三部分 --- 卫星过境。}
黄昏时分，国际空间站
在数分钟内掠过天空。其轨道高度为 \qty{400}{km}，
地球半径为 \qty{6.37e6}{m}，绕行一周耗时 \qty{5.56e3}{s}（约
\qty{93}{min}）。
=\begin{enumerate}[resume]
  \item 在哪一常用参考系中，空间站的运动最为简单？
        求其在该参考系中的速率，单位为 \unit{km/s}。
  \item 用两句话解释，为什么地面参考系描述空间站的运动
        显得别扭 --- 相邻两圈之间，其在地面上的轨迹
        发生了什么变化？
  \item 比较空间站的轨道速率与赤道自转速率
        （$\approx \qty{0.46}{km/s}$），
        空间站快了多少倍？
  \item 一颗地球静止卫星以 \qty{4.22e7}{m} 的轨道半径运行，
        恰好一天绕行一周（$\qty{8.64e4}{s}$）。求其在地心参考系中的速率，
        然后描述其在地面参考系中的运动。
        用一句话协调这两个结果。
  \item 在地心参考系中，月球绕地球运动的轨迹近似为圆，
        周期约为 $27$ 天。试用文字勾勒其在日心参考系中
        几个月内的轨迹。
=\end{enumerate}

=\medskip
=\noindent\textbf{第四部分 --- 谁在运动？}
=\begin{enumerate}[resume]
  \item “太阳从东方升起。”和“地球向东自转。”
        分别说明每句话所描述的观察是在哪一参考系中作出的，
        并解释为何二者都不算错。
  \item 用一段平实的文字回答：托勒密的参考系有哪些优点，
        哥白尼的参考系简化了什么（请提及火星的逆行），
        以及现代的评价为何关乎 \emph{便利性}，
        而不仅仅是某个参考系是否“为真”？
  \item “安乐椅旅行者”的速率报告。你静坐在
        赤道附近的一把安乐椅中。求你在地面参考系中的速率，
        在地心参考系中的速率（$\approx \qty{0.46}{km/s}$），以及
        在日心参考系中的速率（$\approx \qty{30}{km/s}$）。然后计算
        在日心参考系中，你在观看一部
        \qty{90}{min} 的电影期间所经过的距离，并以本章的
        一句话寓意为本题作结。
=\end{enumerate}
=\end{problem}
